\scanpage{114}
Soit \(V\) vectoriel de dim \(2\) sur \(\mathbb F_2\)

\(I\) ens de card. \(3\)

\((P_i)_{i\in I}\) une famille de \(V\)-torseurs

\(\bigwedge_{i\in I}P_i\overset{\varphi}{\simeq}\mathbb{1}\) (torseur trivial) un isom.

\(\Gamma\subset\prod_{i\in I}P_i\) formé des \((x_i)_{i\in I}\) tels que \(\bigwedge_{i\in I}x_i=0\)

(ens. à 16 éléments)

\[
 \begin{aligned}
 \forall i\in I,\quad \widetilde P_i&\xrightarrow{\pi_i}P_i&&\text{un rev. d'ordre 2}\\
 \widetilde \Gamma&\xrightarrow{\pi}\Gamma&&\text{un rev d'ordre 2}
 \end{aligned}\ \bigg|\quad\text{\unsure{autom.|ou torseurs} d'ordre 2 noté }\sigma
\]
\(\forall i\in I\), un isom \(\widetilde \Gamma\simeq \Gamma\times_{P_i}\widetilde P_i\) de rev. de \(\Gamma\) (\deleted{\unsure{section}} associé : \(\widetilde p_i:\widetilde \Gamma\,\unsure{\relbar\joinrel\relbar|\longrightarrow}\,\widetilde P_i\))

Soit \(\widetilde P=\coprod_{i\in I}\widetilde P_i\) (de cardinal \(3\times8=24\))

On introduit dans \(P\) une structure de graphe,

\(\widetilde x,\widetilde y\inserted{\in P\ (x\ne y)}\) liés \unsure{sss|ssi}
\[
 \left|\begin{aligned}
 &\text{ou bien }\exists i\in I\text{ avec }x,y\in\widetilde P_i,\quad y=\sigma x\quad\text{i.e. }\pi_i(x)=\pi_i(y)\\
 &\text{ou bien }\deleted{\exists\widetilde \gamma\in\widetilde \Gamma\ \unsure{\text{et}}}
 x\in\widetilde P_i,\ y\in\widetilde P_j\text{ avec }i\unsure{=|\ne} j,\text{ et }\exists \gamma\in\widetilde \Gamma\text{ avec}\\
 &\hspace{8em}x_{\unsure{1|i}}=\widetilde p_i(\gamma),\ y=\widetilde p_j(\gamma)
 \end{aligned}\right.
\]
\uline{NB} le \(\gamma\) est unique -- \unsure{déf. de lien|déf. d'arêtes|déf. d'étiq.|déjà dit}

\deleted{\uline{Question}}

Soit \(i\ne j\). \deleted{tout \(x\)} \(\forall\,\widetilde x_i\in\widetilde P_i\) et \unsure{tt|\(\forall\)} \(y_j\in P_j\), \(\exists!\,\widetilde y_j\in\widetilde P_j\) sur \(y_j\) lié à \(\widetilde x_i\)

\diagram{\begin{tikzcd}[ampersand replacement=\&,column sep=2em,row sep=1em]
 \widetilde P_i\arrow[d]\&\widetilde P_j\arrow[d]\\P_i\&P_j
\end{tikzcd}}
\uline{Axiome} \(\widetilde P_i\xrightarrow{\alpha_{ji}}\Gamma(\widetilde P_j/P_j)\) est injectif, et \deleted{fait} son image est un torseur sous le sous-groupe \(\unsure{+|t|\Gamma}(P_j)=\operatorname{Ker}(\mathbb F_2^{P_j}\xrightarrow[\text{\unsure{add.} somme}]{}\mathbb F_2)\) de \(\mathbb F_2^{P_j}\) (opérant sur \(\Gamma(\widetilde P_j/P_j)\) qui \deleted{devient}) un torseur \deleted{\unsure{Y}} sous celui-ci)

et les deux injections \(\alpha_{ji},\alpha_{jh}\) \deleted{\unsure{bij.}} de \(\widetilde P_i\) et \(\widetilde P_h\) dans \(\Gamma(\unsure{\overline P_j|\widetilde P_j}/P_j)\) ont des images disjointes
% margin note written upside down
\margin{\unsure{Assemblée} Gén. \unsure{Ex}}
% archive sheet number: 115
